\chapter{Sistema nervoso sensoriale}

% ---------------------------------------------------------------------------
\section{Sensazione e percezione}

La percezione degli stimoli non è la registrazione passiva di un dato fisico: dipende da molti
fattori ed è basata su un'\textbf{interpretazione} dello stimolo, costruita anche sulle
\textbf{esperienze precedenti} e sul \textbf{contesto}. Il caso più immediato da illustrare è
quello visivo.

\rubrica{Esempi di percezione visiva}
\begin{itemize}
  \item \textbf{quadrato inscritto in cerchi concentrici}: i lati del quadrato appaiono curvi,
    mentre in realtà sono perfettamente diritti; rimossi i cerchi, si vedono diritti;
  \item \textbf{figura ambigua ``anatra o coniglio''}: la stessa immagine viene riferita da alcuni
    a un'anatra, da altri a un coniglio; è il contesto in cui la si osserva a decidere;
  \item \textbf{triangolo di Kaniza}: si ``vede'' un triangolo bianco che non è stato disegnato ed
    esiste solo come percezione visiva;
  \item \textbf{dipendenza dal contesto delle dimensioni}: un oggetto posto in secondo piano,
    quando viene ritagliato e portato allo stesso livello dell'oggetto in primo piano, appare
    molto più piccolo, pur essendo in prospettiva delle stesse dimensioni.
\end{itemize}

% ---------------------------------------------------------------------------
\section{Specificità dei sistemi sensoriali}

I sistemi sensoriali sono \textbf{specifici} per ogni tipo di energia a cui l'organismo è
sottoposto: a ciascuna modalità corrispondono un organo recettoriale dedicato, una via e un'area
corticale di destinazione.

\begin{itemize}
  \item \textbf{vista}: l'occhio contiene recettori capaci di interpretare le \textbf{onde
    elettromagnetiche} (luce); il segnale raggiunge la \textbf{corteccia visiva} nel lobo
    occipitale, dove i segnali luminosi non sono solo visti ma \textbf{interpretati};
  \item \textbf{udito}: l'organo recettoriale sta nell'orecchio; i suoni sono interpretati dalla
    \textbf{corteccia uditiva} nel lobo temporale;
  \item \textbf{equilibrio}: nell'orecchio si trova anche il \textbf{sistema vestibolare}, organo
    recettoriale distinto e non correlato al suono; le vie dell'equilibrio proiettano al
    \textbf{cervelletto};
  \item \textbf{gusto e olfatto}: hanno le proprie cortecce dedicate (\textbf{gustativa} e
    \textbf{olfattiva}); le vie olfattive dal naso arrivano direttamente alla corteccia olfattiva;
  \item \textbf{sensi somatici}: provengono da tutto il corpo e sono elaborati nella
    \textbf{corteccia somatosensoriale primaria}, nel lobo parietale, \textbf{dietro la scissura
    di Rolando} (davanti alla scissura sta invece il lobo frontale, sede dei programmi motori).
\end{itemize}

\rubrica{Ruolo del talamo}
La \textbf{maggior parte delle vie sensoriali proietta al talamo}: il talamo elabora le
informazioni e le trasmette ai centri corticali. Fanno eccezione le vie olfattive (dirette alla
corteccia olfattiva) e le vie dell'equilibrio (dirette al cervelletto).

% ---------------------------------------------------------------------------
\section{Compiti dei sistemi sensoriali}

\rubrica{Classificazione per modalità}
\begin{itemize}
  \item \textbf{sistema somatosensoriale}: riceve ed elabora stimoli provenienti dalla superficie
    corporea, dai tessuti profondi e dai visceri; ha diverse modalità percettive, mediate da
    diverse classi di recettori:
    \begin{itemize}
      \item \textbf{tatto};
      \item \textbf{cinestesia} o \textbf{senso articolare}: percezione del proprio corpo, delle
        articolazioni e delle posizioni;
      \item \textbf{temperatura};
      \item \textbf{dolore}.
    \end{itemize}
  \item \textbf{sistema visivo};
  \item \textbf{sistema uditivo};
  \item \textbf{sistema vestibolare};
  \item \textbf{sensi chimici}.
\end{itemize}

\rubrica{Che cosa rilevano}
Tutti questi sistemi rilevano:
\begin{enumerate}
  \item \textbf{cosa} si trova nell'ambiente (esterno e interno);
  \item \textbf{quando} si verificano variazioni;
  \item l'\textbf{entità} della variazione;
  \item la \textbf{sede} della variazione.
\end{enumerate}

% ---------------------------------------------------------------------------
\section{Classificazione per sede dell'informazione}

Le informazioni da elaborare provengono da sedi diverse; su questa base i sistemi sensoriali si
distinguono in:

\begin{itemize}
  \item \textbf{sistemi esterocettivi}: recettori posti sulla superficie corporea, che rilevano le
    variazioni dell'ambiente esterno:
    \begin{itemize}
      \item visivo;
      \item uditivo;
      \item cutaneo;
      \item alcuni sensi chimici (p.\,es.\ gusto e olfatto).
    \end{itemize}
  \item \textbf{sistemi propriocettivi}: rilevano la posizione relativa dei segmenti corporei e
    danno il senso di posizione nello spazio; ne fanno parte i \textbf{fusi neuromuscolari}, gli
    \textbf{organi tendinei del Golgi} e i \textbf{recettori articolari};
  \item \textbf{sistemi enterocettivi}: segnalano eventi corporei interni (p.\,es.\ pressione
    sanguigna, glicemia), che vengono riportati alla norma senza che se ne abbia coscienza; in
    genere \textbf{non raggiungono il livello di coscienza}.
\end{itemize}

\rubrica{I tre momenti comuni a tutte le sensazioni}
\begin{enumerate}
  \item \textbf{stimolo fisico}, di qualsiasi tipo;
  \item insieme di eventi che porta alla \textbf{trasduzione} dello stimolo fisico in impulso
    nervoso (\textbf{messaggio}): p.\,es.\ uno stimolo meccanico come il tatto, o le onde
    elettromagnetiche della luce, vengono convertiti in una differenza di segnale elettrico;
  \item \textbf{risposta al messaggio}.
\end{enumerate}

% ---------------------------------------------------------------------------
\section{Impiego delle informazioni sensoriali}

\rubrica{Informazioni dalla periferia del corpo al SNC}
Sono utilizzate per tre funzioni principali:
\begin{enumerate}
  \item \textbf{percezione della sensazione};
  \item \textbf{controllo del movimento};
  \item \textbf{mantenimento dello stato di veglia}: più si è sottoposti a stimoli, più la
    corteccia cerebrale resta sveglia.
\end{enumerate}

\rubrica{Informazioni dall'interno del corpo al SNC}
Regolano:
\begin{enumerate}
  \item \textbf{temperatura};
  \item \textbf{pressione sanguigna};
  \item \textbf{frequenza cardiaca e respiratoria};
  \item \textbf{movimenti riflessi e volontari}.
\end{enumerate}
Lo stato dell'ambiente interno è quindi in grado di influenzare anche il sistema nervoso somatico
e con esso le performance motorie.

% ---------------------------------------------------------------------------
\section{Le quattro caratteristiche elementari dello stimolo}

I sistemi sensitivi decodificano da uno stimolo quattro caratteristiche elementari.

\rubrica{1. Modalità (qualità)}
Diverse forme di energia vengono trasformate in differenti modalità sensitive dai recettori:
visione, udito, tatto, gusto, olfatto.
\begin{itemize}
  \item visione: onde elettromagnetiche;
  \item udito: onde meccaniche di pressione dell'aria sul timpano;
  \item tatto: stimolazione meccanica dei recettori;
  \item gusto e olfatto: sostanze chimiche. Il gusto presenta anche delle \textbf{submodalità}
    (dolce, salato, amaro, \dots).
\end{itemize}

\rubrica{2. Intensità}
L'intensità o quantità di uno stimolo dipende dallo stimolo stesso. L'intensità di stimolazione
più bassa che si può rilevare è detta \textbf{sensazione soglia}: al di sotto della soglia lo
stimolo non viene più avvertito.
\begin{itemize}
  \item la soglia può variare in seguito a \textbf{utilizzo}, \textbf{fatica} e \textbf{contesto};
  \item ci si può \textbf{adattare} a uno stimolo e non avvertirlo più, oppure al contrario
    \textbf{sensibilizzarsi} e percepirlo anche a intensità molto basse;
  \item \textbf{non cambia la soglia di eccitazione del recettore}, cambia la risposta dei neuroni
    nel SNC, compreso il \textbf{sistema limbico} (quello che si occupa dei comportamenti);
  \item p.\,es.\ un atleta dei 100\,m è molto più pronto nel percepire il colpo di pistola dello
    starter rispetto a chi non ha mai fatto quella gara: un maratoneta è meno rapido di un
    centometrista nel partire al segnale.
\end{itemize}

\rubrica{3. Durata}
È definita dalla relazione tra l'\textbf{intensità dello stimolo} e l'\textbf{intensità
percepita}. Se lo stimolo persiste per lungo tempo, l'intensità percepita diminuisce: il recettore
si abitua allo stimolo secondo il meccanismo dell'\textbf{accomodazione}.

\rubrica{4. Capacità di localizzare}
Capacità di individuare la sede di applicazione dello stimolo e abilità di distinguere la distanza
minima tra due stimoli rilevabili $\rightarrow$ \textbf{soglia dei due punti}.

% ---------------------------------------------------------------------------
\section{I recettori sensoriali}

I \textbf{recettori sensoriali} sono strutture più o meno semplici capaci di rispondere a un solo
tipo di stimolo; non vanno confusi con i recettori di membrana visti nelle lezioni precedenti.

\rubrica{Legge di Müller}
Ogni recettore risponde a un tipo di stimolo. Lo \textbf{stimolo adeguato} è l'unico stimolo che
attiva un recettore specifico e quindi una particolare fibra nervosa.

\rubrica{Morfologia e modalità sensoriali}
A ciascuna modalità corrisponde una morfologia recettoriale diversa:
\begin{itemize}
  \item \textbf{visione};
  \item \textbf{udito} e \textbf{senso dell'equilibrio};
  \item \textbf{gusto};
  \item \textbf{olfatto};
  \item \textbf{propriocezione} della mandibola;
  \item \textbf{tatto, dolore, temperatura, propriocezione} di arti e tronco.
\end{itemize}

\rubrica{Dal più semplice al più complesso}
\begin{itemize}
  \item \textbf{terminazioni nervose libere}: nel neurone pseudounipolare la branca che raggiunge
    la periferia si sfiocca e perde la guaina mielinica; le ramificazioni sono esse stesse i
    recettori sensoriali;
  \item \textbf{corpuscolo di Pacini}: il recettore è la fibra periferica del neurone che ha perso
    la guaina mielinica, ma attorno alla terminazione assonica si è costituita una serie di
    \textbf{lamelle connettivali} con \textbf{liquido interstiziale} all'interno; è impiegato
    p.\,es.\ per la sensazione tattile;
  \item \textbf{bastoncello} (visione): cellula specializzata a trasformare le onde
    elettromagnetiche in potenziali elettrici; non è nemmeno connessa direttamente con la cellula
    nervosa, perché tra recettore e cellula nervosa è interposto un \textbf{interneurone}.
\end{itemize}

% ---------------------------------------------------------------------------
\section{Trasduzione del segnale: il meccanocettore}

Esempio di trasduzione in un meccanocettore (è il funzionamento del \textbf{fuso neuromuscolare},
recettore di lunghezza del muscolo):
\begin{enumerate}
  \item si applica una forza alla cellula: lo \textbf{stiramento} della membrana apre
    \textbf{meccanicamente} il canale ionico;
  \item il $Na^+$ entra nella cellula (flusso ionico in entrata; in uscita il $K^+$);
  \item il potenziale di membrana si sposta da $-70$\,mV verso $-50$/$-40$\,mV fino a raggiungere
    la soglia;
  \item si genera prima il \textbf{potenziale recettoriale}, poi il \textbf{potenziale d'azione}.
\end{enumerate}
L'energia meccanica è stata così trasformata in energia elettrica: è esattamente il compito del
recettore sensoriale.

% ---------------------------------------------------------------------------
\section{Codifica dell'intensità}

Relazione tra intensità dello stimolo, potenziale recettoriale e numero di potenziali d'azione
generati, registrati al \textbf{primo nodo di Ranvier} della fibra afferente:

\begin{itemize}
  \item \textbf{stimolo di bassa intensità}: si aprono i canali per il $Na^+$, il potenziale sale
    p.\,es.\ da $-70$ a $-60$\,mV; \textbf{non si raggiunge la soglia} e alla fine dello stimolo il
    potenziale torna al valore di riposo $\rightarrow$ solo potenziale recettoriale, senza
    conseguenze;
  \item \textbf{intensità maggiore}: potenziale recettoriale più ampio $\rightarrow$ si genera un
    certo numero di potenziali d'azione, che viaggiano dalla periferia verso il midollo spinale
    lungo la fibra afferente;
  \item \textbf{intensità ancora maggiore}: potenziale recettoriale ancora più ampio $\rightarrow$
    aumento del numero di potenziali d'azione.
\end{itemize}

Il \textbf{numero di potenziali d'azione} generati dipende dunque dall'\textbf{intensità dello
stimolo}: la corteccia deputata a decodificare il segnale distingue uno stimolo che ne ha generati
pochi da uno che ne ha generati molti di più, e lo riconosce come più intenso.

\rubrica{Riepilogo}
\begin{itemize}
  \item \textbf{modalità}: dipende dal \textbf{recettore stimolato};
  \item \textbf{intensità}: dipende dal \textbf{numero di potenziali d'azione}.
\end{itemize}

% ---------------------------------------------------------------------------
\section{Adattamento allo stimolo}

In base alla capacità di adattamento (abitudine) allo stimolo si distinguono due classi di
recettori. In entrambi i casi si applica uno stimolo di una certa intensità e durata: si genera un
potenziale recettoriale e, raggiunta la soglia, dei potenziali d'azione.

\rubrica{Recettori a lento adattamento}
\begin{itemize}
  \item il \textbf{potenziale recettoriale} resta pressoché costante, declinando poco per tutta la
    durata di applicazione dello stimolo;
  \item si producono \textbf{potenziali d'azione per tutta la durata dello stimolo}, ma con
    \textbf{frequenza decrescente}: il recettore si sta abituando;
  \item dal decremento della frequenza si ricava \textbf{da quanto tempo lo stimolo è applicato},
    cioè se ne definisce in un certo senso la \textbf{durata};
  \item esempio: i \textbf{fusi neuromuscolari}, che informano continuamente sulla lunghezza del
    muscolo variando la frequenza in base ad accorciamento e allungamento durante il movimento.
\end{itemize}

\rubrica{Recettori a rapido adattamento}
\begin{itemize}
  \item generano il potenziale recettoriale solo \textbf{quando lo stimolo viene applicato}
    (\textit{stimolo on}) e \textbf{quando viene tolto} (\textit{risposta off});
  \item di conseguenza in corteccia arrivano potenziali d'azione solo all'inizio e alla fine:
    si sa \textbf{quando} lo stimolo è applicato e \textbf{quando} è rimosso, ma \textbf{non si
    hanno segnali durante} tutta l'applicazione;
  \item esempio: il \textbf{corpuscolo di Pacini}, recettore molto sensibile. È il recettore che fa
    avvertire un maglione nel momento in cui lo si indossa, e la sensazione di leggerezza quando lo
    si toglie, senza percepirne il peso per tutte le ore in cui lo si porta.
\end{itemize}

\rubrica{Significato energetico}
Una percezione continua sarebbe non solo fastidiosa, ma anche molto dispendiosa: ogni potenziale
d'azione comporta ingresso di $Na^+$ e uscita di $K^+$, e quindi lavoro della \textbf{pompa
sodio-potassio} per ripristinare le concentrazioni, con notevole consumo di ATP e di glucosio. Il
sistema si è perciò evoluto verso recettori a rapido adattamento là dove non è essenziale avere
cognizione continua dello stato di quel particolare recettore o tessuto.
